% Source: https://artofproblemsolving.com/community/c3h2295934_bugging_question
% Stable identity: aops-topic-2295934-q01
\textbf{(a)} For $n=3$: $S(3)=1+2+6=9$, but $\dfrac{3\cdot7\cdot5}{3}=35$. So the conclusion is wrong.

The starting claim fails already for $n=2$: $1!+2!=3$, while $2\cdot1+1\cdot2=4$; for $n=3$ it gives $3\cdot1+2\cdot2+1\cdot3=10\ne9$. In fact
\[\sum_{k=1}^{n}k(n+1-k)=(n+1)\cdot\dfrac{n(n+1)}{2}-\dfrac{n(n+1)(2n+1)}{6}=\dfrac{n(n+1)(n+2)}{6},\]
a cubic in $n$, while $n!$ alone is far larger for large $n$. The claim treats the products $k!=1\cdot2\cdots k$ as if they were sums. (The final formula $\dfrac{n(2n+1)(2n-1)}{3}$ is the sum $1^2+3^2+\cdots+(2n-1)^2$ of the first $n$ odd squares, which has nothing to do with $S(n)$.)

\textbf{(b)} Suppose such a polynomial $P$ exists. Let $Q(x)=P(x)-P(x-1)$, also a polynomial. For every integer $n\ge2$,
\[Q(n)=S(n)-S(n-1)=n!.\]
Hence for every integer $n\ge2$,
\[Q(n+1)=(n+1)!=(n+1)\cdot n!=(n+1)\,Q(n).\]
So the polynomial $Q(x+1)-(x+1)Q(x)$ has infinitely many roots, and is therefore the zero polynomial: $Q(x+1)=(x+1)Q(x)$ for all $x$. Now $Q$ is not the zero polynomial, since $Q(2)=2$. If $Q$ has degree $d$, then $Q(x+1)$ has degree $d$ while $(x+1)Q(x)$ has degree $d+1$, a contradiction. Hence no such polynomial exists. $\blacksquare$
